\section{Исходный код}\label{app:code}

Полный исходный код, конфигурации и сохранённые результаты доступны в
репозитории проекта:
\begin{center}
\url{https://github.com/obilka/fractal_hyperparams_project}
\end{center}
Структура каталогов описана в приложении~\ref{app:repro}.
Ниже приведены ключевые фрагменты.

\subsection{Ансамблевая итерация обучения}

Основной приём реализации --- одновременное обучение всех узлов сетки
гиперпараметров как ансамбля независимых сетей. Веса хранятся с ведущей осью
$G$ (номер конфигурации), гиперпараметры транслируются по осям весов.

\begin{lstlisting}[language=Python,caption={Прямой и обратный проход для ансамбля сетей}]
H  = np.einsum("nd,ghd->gnh", Xb, W0, optimize=True)
Ha = _act(H, cfg.activation)
out = np.einsum("gnh,goh->gno", Ha, W1, optimize=True)
R = out - Yb
loss = np.einsum("gno,gno->g", R, R) / (Nb * d_out)

# ручное обратное распространение
go  = (2.0 / (Nb * d_out)) * R
gW1 = np.einsum("gno,gnh->goh", go, Ha, optimize=True)
gha = np.einsum("gno,goh->gnh", go, W1, optimize=True)
gh  = gha * _act_grad(H, Ha, cfg.activation)
gW0 = np.einsum("gnh,nd->ghd", gh, Xb, optimize=True)

# L2-регуляризация (weight decay)
gW0 = gW0 + wd * W0
gW1 = gW1 + wd * W1
\end{lstlisting}

\subsection{Правила обновления}

\begin{lstlisting}[language=Python,caption={Реализация четырёх оптимизаторов}]
if opt == "sgd":
    W0 -= lr0 * gW0
    W1 -= lr1 * gW1
elif opt == "momentum":
    M0 = beta * M0 + gW0
    M1 = beta * M1 + gW1
    W0 -= lr0 * M0
    W1 -= lr1 * M1
elif opt == "rmsprop":
    V0 = b2 * V0 + (1 - b2) * gW0 * gW0
    V1 = b2 * V1 + (1 - b2) * gW1 * gW1
    W0 -= lr0 * gW0 / (np.sqrt(V0) + eps)
    W1 -= lr1 * gW1 / (np.sqrt(V1) + eps)
elif opt == "adam":
    M0 = b1 * M0 + (1 - b1) * gW0
    V0 = b2 * V0 + (1 - b2) * gW0 * gW0
    bc1, bc2 = 1 - np.power(b1, t + 1), 1 - np.power(b2, t + 1)
    W0 -= lr0 * (M0 / bc1) / (np.sqrt(V0 / bc2) + eps)
    # аналогично для второго слоя
\end{lstlisting}

\subsection{Классификация режимов}

\begin{lstlisting}[language=Python,caption={Классификация траектории на четыре режима}]
m20       = tail[:, -TAIL_SHORT:].mean(axis=1)
m20_first = tail[:, :TAIL_SHORT].mean(axis=1)
rel       = tail.std(axis=1) / np.maximum(tail.mean(axis=1), 1e-300)

stabilized = rel < CONV_TOL              # вышли на стационар
decreasing = m20 < 0.99 * m20_first      # loss продолжает убывать
conv = (m20 < 1.0) & (stabilized | decreasing) & np.isfinite(m20)
div  = (stats["max_l"] > DIV_THRESHOLD) | (m20 > DIV_THRESHOLD)

regime[:]             = OSCILLATING
regime[conv]          = CONVERGED
regime[div & ~nonfin] = DIVERGED
regime[nonfin]        = NONFINITE
\end{lstlisting}

\subsection{Извлечение границы и box-counting}

\begin{lstlisting}[language=Python,caption={Морфологическое извлечение границы и оценка размерности}]
def extract_boundary(regime):
    B = (regime == CONVERGED)
    st = ndimage.generate_binary_structure(2, 1)
    bnd = np.logical_xor(ndimage.binary_dilation(B, structure=st),
                         ndimage.binary_erosion(B, structure=st))
    return skeletonize(bnd) if bnd.sum() > 0 else bnd

def box_counting_dimension(boundary, min_box=2, max_div=8):
    s, sizes, counts = min_box, [], []
    while s <= max(min_box * 2, min(boundary.shape) // max_div):
        ny = (boundary.shape[0] // s) * s
        nx = (boundary.shape[1] // s) * s
        red = boundary[:ny, :nx].reshape(ny // s, s, nx // s, s).any(axis=(1, 3))
        if red.sum() > 0:
            sizes.append(s); counts.append(int(red.sum()))
        s *= 2
    coef, cov = np.polyfit(np.log(sizes), np.log(counts), 1, cov=True)
    return -coef[0], float(np.sqrt(cov[0, 0])), sizes, counts
\end{lstlisting}

\subsection{Теоретический блок}

\begin{lstlisting}[language=Python,caption={Спектральный радиус матрицы перехода метода с инерцией}]
def momentum_spectral_radius(eta_lambda, beta):
    x, b = np.asarray(eta_lambda), np.asarray(beta)
    tr, det = 1.0 + b - x, b
    disc = tr * tr - 4.0 * det
    sq = np.sqrt(np.abs(disc))
    r_real = np.maximum(np.abs((tr + sq) / 2.0), np.abs((tr - sq) / 2.0))
    return np.where(disc >= 0, r_real, np.sqrt(np.abs(det)))
\end{lstlisting}

\section{Инструкция по воспроизведению}\label{app:repro}

\subsection{Структура проекта}

\begin{verbatim}
src/fractal_hp.py        ядро: модель, оптимизаторы, сканирование,
                         box-counting, теоретический блок
src/run_experiments.py   все эксперименты (стадии A-H), resume-режим
src/make_figures.py      построение всех рисунков
src/make_tables.py       генерация LaTeX-таблиц и макросов с числами
report/main.tex          отчёт (XeLaTeX)
report/sections/         разделы отчёта
report/tables/           таблицы, генерируются автоматически
report/figures/          подпись для титульного листа
outputs/data/            результаты сканирований (.npz) и метаданные (.json)
outputs/figures/         рисунки
run_all.sh               полный цикл: расчёты, рисунки, таблицы, сборка PDF
requirements.txt         версии библиотек
\end{verbatim}

\subsection{Требования}

Python 3.10 и выше; NumPy, SciPy, scikit-image, scikit-learn, Matplotlib.
PyTorch не требуется. Для сборки отчёта --- \TeX{}Live с XeLaTeX и пакетами
polyglossia, fontspec, geometry, setspace, booktabs, subcaption, listings,
hyperref; шрифт Liberation Serif. Оперативной памяти достаточно 2 ГБ,
расчёты выполняются на CPU.

\subsection{Запуск}

\begin{lstlisting}[language=bash,caption={Полный цикл воспроизведения}]
git clone https://github.com/obilka/fractal_hyperparams_project.git
cd fractal_hyperparams_project
pip install -r requirements.txt
bash run_all.sh
\end{lstlisting}

Скрипт последовательно выполняет расчёты (стадии A--H), строит рисунки и
таблицы и собирает PDF. Отдельные стадии запускаются как
\texttt{python3 src/run\_experiments.py B C}. Расчёт поддерживает
возобновление: результат каждого блока из 4096 конфигураций сохраняется в
файл контрольной точки, и повторный запуск продолжает работу с места
остановки, а уже завершённые эксперименты пропускаются.

Суммарное машинное время всех экспериментов --- \TotalHours~ч на одном ядре
CPU.

\subsection{Фиксация случайности и точности}

Seed равен 42 для генерации данных, $42+1000$ для инициализации весов и
$42+7$ для порядка примеров в мини-батчах. Инициализация одинакова для всех
узлов сетки в пределах одного эксперимента. Тип данных по умолчанию
\texttt{float64}; эксперимент G1 выполнен в \texttt{float32} для оценки
влияния точности. Все параметры запуска сохраняются в JSON рядом с
результатом.

\section{Дополнительные материалы}

Числовые значения, приведённые в тексте и таблицах отчёта, подставляются
автоматически из файлов \texttt{outputs/data/*.json} через генерируемый файл
\texttt{report/tables/numbers.tex}. Это гарантирует, что текст отчёта,
таблицы и сохранённые результаты расчётов не могут разойтись между собой:
при повторном запуске расчётов и пересборке отчёта все числа обновляются
согласованно.
